% =====================================================================
% Appendix B --- The closure first-order conditions
% ---------------------------------------------------------------------
% NEW (2026-10-15). Absorbs the condensed closure prose of Section 4.3
% ("Five labour closures and three financings"), which leaves Section 4 as
% the kernel only. Job: state, in one place and compactly, the condition
% that defines each labour-market closure and each financing closure, and
% fix the notation the matrix tables read their symbols off.
% The functional forms and the equations themselves live in the model
% appendix (Appendix~\ref{app:model}); the conditions are \eqref'd there, so
% this appendix carries NO duplicate equation label. Equivalences between
% the closures are in Appendix~\ref{app:equivalences}.
% =====================================================================
\section{The closure first-order conditions}\label{app:closure-focs}

\noindent
Each closure adds one condition to the kernel of
Section~\ref{sec:kernel}. This appendix states those conditions --- the
five labour-market closures first, then the three financing closures --- and
fixes the notation the matrix tables use. The functional forms and the
equations themselves are in Appendix~\ref{app:model}; the equivalences
between the closures are in Appendix~\ref{app:equivalences}.

\subsection{Notation}\label{app:foc-notation}

$L_i$ is employment in sector $i$ and $L=\sum_i L_i$ total employment; a bar
marks the baseline, so $\bar L_i$ is baseline sectoral and $\Lbar$ baseline
total employment. $w$ is the common wage of the mobile closures, $w_i$ the
sector-specific wage where labour markets are sectoral, and $\bar w$ the
pegged real wage. Real wages are deflated by the consumption price index
$\Pi(p)$, the CES price of the household's bundle; a superscript $0$ marks
baseline values. $L_i^{\mathrm{cm}}(p,y_i,w_i)$ is the labour demand that
solves the firm's cost-minimization problem
of Section~\ref{subsec:optimization} at prices $p$ and output $y_i$. The
allocation parameter $\eta$ takes only its endpoints: $\eta=1$ lets labour
move to its cost-minimizing use, $\eta=0$ freezes the baseline allocation.
The labour-supply elasticity $\eta_{s,i}$ governs how strongly sectoral
employment responds to the real wage; the headline results set it to a
common value $\eta_s$. Finally, $\theta$, $\epsilon$ and $\sigma$ are the
substitution elasticities of Section~\ref{subsec:optimization} (among
intermediates, between labour and intermediates, and in consumption); their
limit $0^{+}$ is fixed proportions.

On the financing side, $c_i^h$ is the household purchase of good $i$ and
$E^h$ its consumption expenditure, so $\sum_i p_i c_i^h=E^h$ is the
household budget of~\eqref{eq:utility-max}, held fixed while F1 tilts its
composition; $W=\sum_i w_i L_i$ is the wage bill, $T_G(p)=\sum_i p_i g_i^G$
the baseline tax cost of government purchases, and $g_i$ the programme
bundle (real public investment in sector $i$) with current-price cost
$\sum_i p_i g_i$. Under F2 that cost is raised as the programme tax $T(p)$;
under F3 it is booked as external programme financing $B_{\mathrm{gov}}$.

\subsection{The five labour-market conditions}\label{app:foc-labour}

Each condition below is imposed on the common labour
demand~$L_i^{\mathrm{cm}}$; the closures differ only in the wage and
labour-supply restriction they add.

\begin{itemize}
\item \textbf{The no-reallocation closure.} Employment is fixed in every
  sector and a sector-specific wage clears each labour market,
  $L_i=L_i^{\mathrm{cm}}(p,y_i,w_i)=\bar L_i$, $i=1,\ldots,N$, $\eta=0$
  (eq.~\ref{eq:closure-bf}).
\item \textbf{The standard CGE closure.} Labour reallocates across sectors
  at one flexible wage while total employment stays fixed,
  $\sum_i L_i=\Lbar$, $\eta=1$ (eq.~\ref{eq:closure-alpha}).
\item \textbf{The elastic labour-supply closure.} Each sector has its own
  wage and a supply curve on the real wage with a common elasticity,
  $L_i=\bar L_i[(w_i/\Pi(p))/(w_{i0}/\Pi_0)]^{\eta_s}$, $\eta_s\geq0$
  (eq.~\ref{eq:closure-beta}); its $N$ markets determine the $N$ wages.
\item \textbf{The fixed real wage closure.} The real wage is pegged at its
  baseline and employment is endogenous and uncapped,
  $w/\Pi(p)=\bar w=1$, $\eta=1$ (eq.~\ref{eq:closure-gamma}).
\item \textbf{The Leontief closure.} The fixed real wage rule combined with
  the zero-substitution limit,
  $(\theta,\epsilon,\sigma)\to(0^{+},0^{+},0^{+})$
  (eq.~\ref{eq:closure-delta}); it is a corner of the fixed real wage
  closure, not a separate labour-market mechanism.
\end{itemize}

\subsection{The three financing conditions}\label{app:foc-financing}

Each labour closure is crossed with one of three ways of accommodating the
same programme scale and sectoral incidence. Baseline government purchases
$g^G$ remain tax-financed in all three cases; let $R$ be the net external
transfer entering household resources after tax ($R$ is the variable called
$F$ in the kernel and the run manifests).

\begin{itemize}
\item \textbf{F1: budget-neutral preference reallocation.} The programme
  scale and incidence tilt the household's CES weights without adding a
  public-demand bundle, $\widetilde\omega_i=\omega_i d_i/\sum_j\omega_j d_j$,
  $E=W-T_G(p)+R$ (eq.~\ref{eq:closure-f1}). Spending shares rise in
  programme sectors and fall elsewhere within the household budget.
\item \textbf{F2: tax-financed public investment.} The programme bundle $g$
  is added and financed by an additional household tax equal to its
  current-price cost, $T(p)=\sum_i p_i g_i$, $E=W-T_G(p)-T(p)+R$,
  $B_{\mathrm{gov}}=0$ (eq.~\ref{eq:closure-f2}).
\item \textbf{F3: externally financed public investment.} The same bundle,
  with no additional domestic programme tax; its cost is booked as external
  programme financing, $B_{\mathrm{gov}}=\sum_i p_i g_i$,
  $E=W-T_G(p)+R$ (eq.~\ref{eq:closure-f3}).
\end{itemize}

Under every combination all $N$ goods markets clear
(eq.~\ref{eq:model-clearing}) and the external-account identity holds
(eq.~\ref{eq:financing-external-account}). The $R$ variable is solved
endogenously in the no-reallocation, standard CGE and elastic-supply
closures, but fixed at zero in the two fixed-wage closures.
